\chapter{Inference for Two Proportions}\index{two-proportion inference}

Does a text reminder make students more likely to return a permission slip? Do
juniors get enough sleep more often than seniors? Does a new vaccine prevent
infection better than a placebo? Each question compares two groups, and each answer
is a proportion: a fraction of ``yes'' in one group against a fraction of ``yes'' in
the other.

The Statistical Inferences for Proportions chapter studied one proportion at a time,
and the Inference for Two Means chapter compared two means. This chapter puts
the two ideas together: confidence intervals and significance tests for the
\emph{difference} between two population proportions. The reasoning is the same as
before. The only new piece is figuring out how much a difference varies from sample
to sample.

\section{The Difference Between Two Proportions}
Suppose a school counselor wants to compare how many juniors and seniors get at least 8
hours of sleep on school nights. Let $p_1$ be the true proportion of juniors and $p_2$ the
true proportion of seniors. The parameter of interest is the difference
$$p_1 - p_2$$
The counselor surveys random samples: 84 of 240 juniors ($\hat{p}_1 = 0.35$) and 54 of 200
seniors ($\hat{p}_2 = 0.27$) say they sleep 8 hours or more. The point estimate is
$\hat{p}_1 - \hat{p}_2 = 0.08$: juniors are 8 percentage points more likely to get enough sleep
\emph{in these samples}. Would another pair of samples show the same gap? To know how much
to trust 0.08, we need the sampling distribution of $\hat{p}_1 - \hat{p}_2$.

\subsection{Variances Add, Even When You Subtract}\index{variance!of a difference}
From the Sampling Distributions and the Central Limit Theorem chapter, each sample
proportion has mean $p$ and standard deviation $\sqrt{\frac{p(1-p)}{n}}$. The Random
Variables chapter showed that for independent random variables, variances add. That is
true for differences as well as sums: $\operatorname{Var}(X - Y) = \operatorname{Var}(X) + \operatorname{Var}(Y)$,
because subtracting a random quantity adds just as much uncertainty as adding it. So:

\begin{mdframed}[style = important, frametitle = {Sampling Distribution of $\hat{p}_1 - \hat{p}_2$}]
For independent random samples,
$$\mu_{\hat{p}_1 - \hat{p}_2} = p_1 - p_2, \qquad
\sigma_{\hat{p}_1 - \hat{p}_2} = \sqrt{\frac{p_1(1-p_1)}{n_1} + \frac{p_2(1-p_2)}{n_2}}$$
The distribution is approximately normal when the large-counts condition holds in both groups.
\end{mdframed}

It is tempting to subtract the two standard deviations instead. A simulation shows how
wrong that would be. Here are 10,000 pairs of surveys with $p_1 = 0.35$ and $p_2 = 0.27$:
\begin{minted}{python}
import numpy as np

rng = np.random.default_rng(3)
p1, n1 = 0.35, 240        # juniors
p2, n2 = 0.27, 200        # seniors

# 10,000 pairs of surveys
p_hat_1 = rng.binomial(n1, p1, size=10_000) / n1
p_hat_2 = rng.binomial(n2, p2, size=10_000) / n2
diff = p_hat_1 - p_hat_2

theory = np.sqrt(p1 * (1 - p1) / n1 + p2 * (1 - p2) / n2)
print(f"mean of differences: {diff.mean():.4f}  (p1 - p2 = {p1 - p2:.2f})")
print(f"sd of differences:   {diff.std():.4f}  (formula: {theory:.4f})")
print(f"sd if you subtracted the sds: {abs(np.sqrt(p1*(1-p1)/n1) - np.sqrt(p2*(1-p2)/n2)):.4f}")
\end{minted}
\begin{verbatim}
mean of differences: 0.0794  (p1 - p2 = 0.08)
sd of differences:   0.0439  (formula: 0.0440)
sd if you subtracted the sds: 0.0006
\end{verbatim}
The differences really do spread out by about 0.044, just as the formula says. Subtracting
the standard deviations would claim the difference is almost perfectly predictable, which is
nonsense: two uncertain numbers give a \emph{more} uncertain difference, not a less uncertain one.

\section{Conditions}\index{conditions for inference!two proportions}
Before any two-proportion interval or test, check:
\begin{itemize}
\item \textbf{Random:} the data come from two independent random samples, or from two groups
in a randomized experiment.
\item \textbf{10\%:} when sampling without replacement, each sample is no more than 10\% of its
population. (This is not needed for a randomized experiment.)
\item \textbf{Large counts:} for a confidence interval, the four counts $n_1\hat{p}_1$,
$n_1(1-\hat{p}_1)$, $n_2\hat{p}_2$, and $n_2(1-\hat{p}_2)$ are all at least 10. In other words,
each group has at least 10 ``yes'' answers and at least 10 ``no'' answers.
\end{itemize}

\section{A Confidence Interval for $p_1 - p_2$}\index{confidence interval!for a difference in proportions}
We do not know $p_1$ and $p_2$, so we estimate the standard deviation with the sample
proportions. The result is the \newterm{standard error} of the difference, and the interval
has the familiar shape: estimate $\pm$ critical value $\times$ standard error.

\begin{mdframed}[style = important, frametitle = {Two-Proportion $z$-Interval}]
$$(\hat{p}_1 - \hat{p}_2) \pm z^*\sqrt{\frac{\hat{p}_1(1-\hat{p}_1)}{n_1} + \frac{\hat{p}_2(1-\hat{p}_2)}{n_2}}$$
\end{mdframed}

\begin{mdframed}[frametitle = {Example}, style = important]
For the sleep survey, the counts are 84 and 156 for juniors and 54 and 146 for seniors, all
at least 10, and the samples are random and less than 10\% of each class in a large district.
The standard error is
$$SE = \sqrt{\frac{0.35(0.65)}{240} + \frac{0.27(0.73)}{200}} \approx 0.0440$$
For 95\% confidence, $z^* = 1.96$, so the margin of error is $1.96(0.0440) \approx 0.086$, and the
interval is
$$0.08 \pm 0.086, \quad\text{or}\quad (-0.006,\ 0.166)$$
We are 95\% confident that the proportion of juniors who sleep at least 8 hours is between 0.6
percentage points \emph{lower} and 16.6 points \emph{higher} than the proportion of seniors.
\end{mdframed}

Look at that interval closely. It contains 0. So although the samples showed an 8-point gap, the
data are consistent with juniors and seniors sleeping equally well. An interval for a
difference that contains 0 does not give convincing evidence of any difference at all; one
that lies entirely above (or below) 0 does.

\section{A Significance Test for $p_1 - p_2$}\index{significance test!two proportions}
A school wants to know whether texting families a reminder raises the rate of returned
permission slips. In a randomized experiment, 150 families get a reminder text and 150 do not.
In the text group, 102 return the slip ($\hat{p}_1 = 0.68$); in the other group, 84 do
($\hat{p}_2 = 0.56$). Is a 12-point gap convincing?

The hypotheses are
$$H_0: p_1 - p_2 = 0 \qquad H_a: p_1 - p_2 > 0$$
where $p_1$ and $p_2$ are the true return rates with and without a reminder.

\subsection{Pooling}\index{pooled proportion}
A significance test always asks, ``How surprising is this result \emph{if $H_0$ is true}?'' If
$H_0$ is true, the two groups share a single return rate. Our best estimate of that common rate
combines both samples into one \newterm{pooled proportion}:
$$\hat{p}_c = \frac{\text{total successes}}{\text{total sample size}} = \frac{102 + 84}{150 + 150} = \frac{186}{300} = 0.62$$
In the test, we use $\hat{p}_c$ in place of both $\hat{p}_1$ and $\hat{p}_2$ in the standard error,
and in the large-counts check: $150(0.62) = 93$ and $150(0.38) = 57$, both at least 10.

\begin{mdframed}[style = important, frametitle = {Two-Proportion $z$-Test}]
$$z = \frac{(\hat{p}_1 - \hat{p}_2) - 0}{\sqrt{\hat{p}_c(1-\hat{p}_c)\left(\frac{1}{n_1} + \frac{1}{n_2}\right)}}$$
\end{mdframed}

Pool for a test (because $H_0$ says the proportions are equal), but not for an interval (because an
interval makes no such assumption). This is the single most common mistake on two-proportion
problems.

\begin{mdframed}[frametitle = {Example}, style = important]
$$SE = \sqrt{0.62(0.38)\left(\frac{1}{150} + \frac{1}{150}\right)} \approx 0.0560, \qquad
z = \frac{0.68 - 0.56}{0.0560} \approx 2.14$$
The p-value is $P(Z \geq 2.14) \approx 0.016$. Since $0.016 < 0.05$, we reject $H_0$. There is
convincing evidence that reminder texts raise the return rate.
\end{mdframed}

Because families were \emph{randomly assigned} to get a text or not, the two groups differed only
by chance and by the text itself. So we can conclude that the text \emph{caused} the higher return
rate. In the sleep survey, which was an observational study, even a significant difference
would not show that being a junior causes better sleep. (The Experimental Design chapter explains
why random assignment makes this difference.)

\section{Computing}
\textbf{Excel.} With $\hat{p}_1$, $\hat{p}_2$, $n_1$, and $n_2$ in cells, the standard error and
interval are ordinary formulas, and the p-value of a one-sided test is
\texttt{=1-NORM.S.DIST(z, TRUE)}. For the test above, \texttt{=1-NORM.S.DIST(2.14, TRUE)} gives
about 0.016.

\textbf{Python.} The \texttt{statsmodels} library does both procedures from the counts:
\begin{minted}{python}
from statsmodels.stats.proportion import proportions_ztest, confint_proportions_2indep

# Confidence interval: juniors vs. seniors who sleep 8+ hours
low, high = confint_proportions_2indep(84, 240, 54, 200, method="wald")
print(f"95% CI for p1 - p2: ({low:.3f}, {high:.3f})")

# Significance test: text reminder vs. no reminder
z, p_value = proportions_ztest([102, 84], [150, 150], alternative="larger")
print(f"z = {z:.2f}, p-value = {p_value:.4f}")
\end{minted}
\begin{verbatim}
95% CI for p1 - p2: (-0.006, 0.166)
z = 2.14, p-value = 0.0161
\end{verbatim}
The option \mintinline{python}{method="wald"} gives the textbook interval from this chapter;
\texttt{statsmodels} also offers more accurate variations.

\begin{Exercise}[title={Biking to School}, label=tp_ci]
At School A, a random sample shows that 45 of 150 students bike to school. At School B, 30 of 160
do. Construct and interpret a 90\% confidence interval for the difference in the proportions of
students who bike. Is there convincing evidence of a difference?
\end{Exercise}
\begin{Answer}[ref=tp_ci]
$\hat{p}_1 = 0.30$ and $\hat{p}_2 = 0.1875$. The counts 45, 105, 30, and 130 are all at least 10.
$$SE = \sqrt{\frac{0.30(0.70)}{150} + \frac{0.1875(0.8125)}{160}} \approx 0.0485$$
With $z^* = 1.645$, the margin of error is about $0.080$, and the interval is
$0.1125 \pm 0.080$, or about $(0.033, 0.192)$. We are 90\% confident that the proportion of
students who bike is between 3.3 and 19.2 percentage points higher at School A. The entire
interval is above 0, so there is convincing evidence of a difference.
\end{Answer}

\begin{Exercise}[title={A New Math App}, label=tp_test]
In a randomized experiment, 200 students used a new math app and 200 used the old one. On the
unit test, 132 of the new-app students passed and 110 of the old-app students passed. Test, at
$\alpha = 0.05$, whether the pass rates differ (a two-sided test).
\end{Exercise}
\begin{Answer}[ref=tp_test]
$H_0: p_1 - p_2 = 0$ and $H_a: p_1 - p_2 \neq 0$. The pooled proportion is
$\hat{p}_c = \frac{242}{400} = 0.605$, and the pooled counts $200(0.605) = 121$ and $200(0.395) = 79$
are at least 10.
$$SE = \sqrt{0.605(0.395)\left(\frac{1}{200} + \frac{1}{200}\right)} \approx 0.0489, \qquad
z = \frac{0.66 - 0.55}{0.0489} \approx 2.25$$
The two-sided p-value is $2P(Z \geq 2.25) \approx 0.024$. Since $0.024 < 0.05$, reject $H_0$: there is
convincing evidence that the pass rates differ, and because this was a randomized experiment,
the new app caused the difference.
\end{Answer}

\begin{Exercise}[title={Reading an Interval}, label=tp_read]
A 95\% confidence interval for $p_1 - p_2$, the difference in the proportions of adults and teens
who prefer a new soda, is $(-0.02, 0.10)$. A news site reports, ``Adults prefer the new soda more
than teens do.'' Is that justified?
\end{Exercise}
\begin{Answer}[ref=tp_read]
No. The interval includes 0, and even some negative values, so it is plausible that the two
proportions are equal, or that teens prefer it slightly more. The data do not give convincing
evidence of a difference.
\end{Answer}

\begin{Exercise}[title={To Pool or Not to Pool?}, label=tp_pool]
Explain why a two-proportion significance test of $H_0: p_1 = p_2$ uses the pooled proportion
$\hat{p}_c$, but a confidence interval for $p_1 - p_2$ does not.
\end{Exercise}
\begin{Answer}[ref=tp_pool]
A test computes how surprising the data would be \emph{if $H_0$ were true}, and $H_0$ says the two
proportions are the same, so the best estimate of that common value combines both samples. A
confidence interval makes no assumption that the proportions are equal (in fact, it is estimating
how different they are), so it uses each sample's own proportion.
\end{Answer}

\begin{Exercise}[title={What Could Go Wrong?}, label=tp_errors]
In the reminder-text experiment, the school rejected $H_0$. Describe a Type I error and a Type II
error in this context. Which one could the school have made?
\end{Exercise}
\begin{Answer}[ref=tp_errors]
A Type I error would be concluding that reminder texts raise the return rate when they actually
make no difference. A Type II error would be failing to find that the texts help when they really
do. Since the school rejected $H_0$, the only error it could have made is a Type I error.
\end{Answer}
